\documentclass[11pt]{article}

\usepackage[margin=1in]{geometry}
\usepackage{amsmath}
\usepackage{amssymb}
\usepackage{booktabs}
\usepackage{longtable}
\usepackage{caption}
\usepackage{graphicx}
\usepackage[babel=false]{microtype}
\usepackage{xurl}
\usepackage{url}
\usepackage[hidelinks]{hyperref}

\hbadness=10000
\vbadness=10000

\renewcommand{\thetable}{S\arabic{table}}
\renewcommand{\thefigure}{S\arabic{figure}}

\title{Supporting Information\\[0.5em]\large COMPASS: Retrospective Benchmarking of Transferred Slow-Coordinate Priors (BTP report)}
\author{Anmol Goel and Ishaan Gupta\\Department of Biochemical Engineering and Biotechnology\\Indian Institute of Technology Delhi}
\date{}

\begin{document}
\maketitle

\tableofcontents
\newpage

\section{Molecular dynamics protocol (source and target trajectories)}

Production trajectories for the wild-type alanine-rich peptide A10 (Ace-(Ala)$_{10}$-NMe) and the mutant A8Q2 (Ace-(Ala)$_8$(Gln)$_2$-NMe) were generated with \textsc{OpenMM}~8.x.\cite{eastman2017openmm}
We used the additive CHARMM36m protein force field with CHARMM-modified TIP3P water as distributed with \textsc{OpenMM} (\path{charmm36_2024.xml} and matching water parameters).\cite{huang2017charmm36m,kim2015charmmgui}
Nonbonded interactions used a nominal 1.2~nm cutoff with smoothing from 1.0~nm, particle-mesh Ewald electrostatics, constraints on bonds to hydrogen, and rigid water.
Integration used a Langevin middle integrator with 2~fs timestep and friction 1~ps$^{-1}$.
Each system was energy-minimized, equilibrated (NVT with peptide heavy-atom restraints followed by NPT), and simulated in the NVT ensemble for production at \textbf{330~K} to accelerate helix--coil sampling relative to room temperature while preserving the qualitative physics relevant to this benchmark.
The same protocol and temperature were used for A10 and A8Q2 so relative comparisons are internally consistent.

Peptides were solvated in a periodic cubic box with roughly 1.2--1.5~nm buffer and 0.15~M NaCl (see project \path{docs/EXPERIMENT_PROTOCOL.md} for box and ion counts).
Production coordinates were saved at \textbf{0.1~ns} (100~ps) spacing for downstream featurization, consistent with the chunk library described in the main text.
The full A8Q2 trajectory was partitioned into non-overlapping chunks for COMPASS replay; reference observables $f_{\mathrm{helix}}^{\mathrm{ref}}$ and $f_{\mathrm{open}}^{\mathrm{ref}}$ were computed from the same full trajectory.

\section{TICA and SPIB training details}

TICA was fit on backbone $\phi/\psi$ (and related) torsional features from the A10 source trajectory using PyEMMA's symmetric time-lagged covariance estimator at a lag matched to the implementation used for production analysis (order 1--4~ns in frame units; see repository \path{src/} and phase-5 analysis logs).
SPIB encoders were trained in PyTorch on the same A10 feature pipeline as in the main experiments, with latent dimension one or two and information-bottleneck trade-off as reported in the main text.

\section{Replay implementation}

Chunk scoring, deficit modes (raw vs.\ information-gain), two-phase schedules, prior corruption (uniform, shuffled, inverted, noisy), and bootstrap random baselines ($100$ repeats) were implemented in Python.
Code paths: \path{src/compass/replay_enhanced.py} and \path{scripts/compass_experiments_enhanced.py}.
Random seeds and experiment identifiers for each row of Table~\ref{tab:si_sweep} are recorded under \path{archive/outputs_legacy/compass_sweep2/sweep/}.

\section{Complete sweep metrics}

Table~\ref{tab:si_sweep} lists final metrics for every experiment in the enhanced sweep (\path{sweep2_results.csv}).
Columns are: number of acquisition steps $N_{\mathrm{steps}}$, total acquired time $T$, TICA- or SPIB-space JSD for COMPASS vs.\ random mean, and helix-fraction error for COMPASS vs.\ random mean.

\bgroup\scriptsize\setlength{\tabcolsep}{2.5pt}
\begin{longtable}{@{}p{0.52\linewidth}rrrrrr@{}}
\caption{Complete retrospective sweep (enhanced COMPASS, \texttt{sweep2}). Subscripts C/R denote COMPASS and bootstrap random mean ($100$ repeats).}\label{tab:si_sweep}\\
\toprule
Experiment & $N_{\mathrm{steps}}$ & $T$ (ns) & JSD$_\mathrm{C}$ & JSD$_\mathrm{R}$ & $\epsilon_{\mathrm{helix,C}}$ & $\epsilon_{\mathrm{helix,R}}$ \\
\midrule
\endfirsthead
\multicolumn{7}{c}{\tablename\ \thetable{} --- \textit{continued from previous page}} \\
\toprule
Experiment & $N_{\mathrm{steps}}$ & $T$ (ns) & JSD$_\mathrm{C}$ & JSD$_\mathrm{R}$ & $\epsilon_{\mathrm{helix,C}}$ & $\epsilon_{\mathrm{helix,R}}$ \\
\midrule
\endhead
\midrule
\multicolumn{7}{r}{\textit{Continued on next page}} \\
\endfoot
\bottomrule
\endlastfoot
e01\_tica2d\_static\_raw\_80ns & 80 & 80 & 0.1289 & 0.1315 & 0.0105 & 0.01708 \\
e02\_tica2d\_static\_ig\_80ns & 80 & 80 & 0.1233 & 0.1315 & 0.00925 & 0.01708 \\
e03\_tica2d\_two\_phase\_ig\_80ns & 80 & 80 & 0.1473 & 0.1315 & 0.0045 & 0.01708 \\
e04\_tica2d\_anneal\_ig\_80ns & 80 & 80 & 0.1401 & 0.1315 & 0.01175 & 0.01708 \\
e05\_tica1d\_two\_phase\_ig\_80ns & 80 & 80 & 0.0164 & 0.0145 & 0.0545 & 0.01708 \\
e06\_spib1d\_two\_phase\_ig\_80ns & 80 & 80 & 0.0141 & 0.009251 & 0.08325 & 0.01708 \\
e07\_spib2d\_two\_phase\_ig\_80ns & 80 & 80 & 0.1012 & 0.09976 & 0.0895 & 0.01708 \\
e08\_tica2d\_two\_phase\_ig\_200ns & 200 & 200 & 0.04201 & 0.05742 & 0.024 & 0.01097 \\
e09\_spib1d\_two\_phase\_ig\_200ns & 200 & 200 & 0.008417 & 0.003508 & 0.0415 & 0.01097 \\
e10\_spib2d\_two\_phase\_ig\_200ns & 200 & 200 & 0.03434 & 0.04299 & 0.0315 & 0.01097 \\
e11\_tica2d\_two\_phase\_chunk5\_80ns & 160 & 80 & 0.141 & 0.125 & 0.0245 & 0.01181 \\
e12\_tica2d\_two\_phase\_chunk20\_80ns & 40 & 80 & 0.1556 & 0.1386 & 0.00075 & 0.0208 \\
e13\_tica2d\_uniform\_prior\_80ns & 80 & 80 & 0.1997 & 0.1315 & 0.0695 & 0.01708 \\
e14\_tica2d\_shuffled\_prior\_80ns & 80 & 80 & 0.1507 & 0.1315 & 0.02075 & 0.01708 \\
e15\_tica2d\_inverted\_prior\_80ns & 80 & 80 & 0.273 & 0.1315 & 0.0645 & 0.01708 \\
e16\_tica2d\_noisy\_prior\_80ns & 80 & 80 & 0.1403 & 0.1315 & 0.0245 & 0.01708 \\
e17\_tica2d\_pure\_frontier\_schedule\_80ns & 80 & 80 & 0.1997 & 0.1315 & 0.0695 & 0.01708 \\
e18\_tica2d\_exploratory\_static\_80ns & 80 & 80 & 0.216 & 0.1315 & 0.0445 & 0.01708 \\
e19\_tica2d\_two\_phase\_softmax\_repeats\_80ns & 80 & 80 & 0.1442 & 0.1315 & 0.01198 & 0.01708 \\
e20\_spib1d\_two\_phase\_softmax\_repeats\_80ns & 80 & 80 & 0.01721 & 0.009251 & 0.07544 & 0.01708 \\
e21\_tica2d\_chunk20\_softmax\_confirm\_80ns & 40 & 80 & 0.1572 & 0.139 & 0.01172 & 0.02372 \\
\end{longtable}

\egroup

\section{Auxiliary experiments (referenced from the main text)}

Stochastic top-$k$ (softmax) confirmation runs, pure-frontier and exploratory-static ablations, and additional chunk-size variants are included in Table~\ref{tab:si_sweep} (experiments \texttt{e11}--\texttt{e20}).
Score-decomposition plots and per-step decision logs for each run are under \path{archive/outputs_legacy/compass_sweep2/} in the project tree.

\bibliographystyle{unsrt}
\bibliography{COMPASS_JCTC_REFERENCES}

\end{document}
